\HMTNarrativeHeading{What the constant couples}

The holographic arithmetic-geometric coupling constant is introduced at the
point where its register is produced. Its integral representation is denoted
by \(K\); its periodic reading, by \(\kappa_{\rm ag}\). The phase
origin, orientation, and order of the components form part of the definition.
The word ``constant'' names the object determined by this construction and
by its readers, which will be presented through explicit operations.

The extraction brings together the oriented channels and a sum datum.
Hadamard inversion reconstructs their components in the natural order.
Each coefficient of that inversion comes from the declared operator and
its integral image. In this way, the register is determined before it is
used in the closure of the fine-structure constant. The positional reading
then converts the twelve components into a periodic rational expression.
Once the base and length are fixed, this conversion admits an exact inverse
that restores all the blocks.

The same ordering has a geometrical function. Its readers distinguish a
marked tetrad and an incidence configuration; the signature associated with
alpha selects a hexad within the corresponding design. The relation to the
numerical regions uses the alignment and lifting maps that will be proved
in the exceptional part. The periodic reading and the incidence
configuration thus receive a common source, through operations of different
types.

The subsequent operator development adds a third function. The cyclic shifts
of the register form a basis of its carrier. Responses on this basis
determine an operator, and the Gram form provides a stability bound for its
recovery. When the specified cyclic symmetry is present, one complete
vector-valued response makes it possible to reconstruct the others. The
isometric archive preserves this capacity for identification.

This is the content of the coupling that structures the exposition:
reversible arithmetic encoding, marked incidence, and operator
identification. The register is constructed first; what each reader obtains
from it is then proved. Recovery theory will be developed once its spaces
and transports are available. This position explains both the centrality
of the constant and the hierarchy of the proofs that use its different
realizations.
